\documentclass[12pt,a4paper]{article} % Maybe use 'report' ???
% REMEMBER TO REMOVE THE 'draft' TAG!
\usepackage{polyglossia}	% XeTeX multi-lingual support
\usepackage{fontspec}		% The XeTeX font spec package
\usepackage{xunicode}		% XeTeX Unicode character support
\usepackage{xltxtra}		% Umm something else XeTeX

\usepackage[intlimits]{amsmath}	% Math stuff
\usepackage{listings}	% Package for code formatting
\usepackage[dvipsnames]{xcolor}
%\include{/home/johnlm/code/avr_assembler_listing_def.tex}
%\input{avr_assembler_listing_def.tex}
%\renewcommand{\lstlistingname}{Koda bloks}
\usepackage{booktabs}	% Package for nicer tables
\usepackage[perpage,symbol*]{footmisc}	% Resets footnote marks per page
%\usepackage{wrapfig}
\usepackage[superscript]{cite}	% Superscript cites
\usepackage{rotating}	% Package for landscape figures
\usepackage{indentfirst}	%Indent first paragraph too
\usepackage{setspace}	% Line spacing package
\usepackage{fancyhdr}	% Fancy header/footer package
\usepackage[figurewithin=section,tablewithin=section]%
	{caption}	% Package for configuring captions format
\usepackage[a4paper,top=20mm, bottom=20mm, left=35mm, right=20mm]%
%\usepackage[a4paper,top=22mm, bottom=25mm, left=35mm, right=20mm]% DEBUG
	{geometry}
\usepackage{fixlatvian}	% Fixes for Latvian typography rules

% Preamble formatting settings defined included from external file
% Polyglossia settings
\setmainlanguage{latvian}
\setotherlanguages{english,russian}

% FONTS (Empty \setmainfont gives you new Computer Modern)
% Deja Vu Family (rather good but wide chars)
%\setmainfont{DejaVu Serif}
%\setsansfont{DejaVu Sans}

% Liberation Family (my current favourite)
%\setmainfont{Liberation Serif}
%\setsansfont{Liberation Sans}

% Monos
%\setmonofont{DejaVu Sans Mono}

% Russian Font
%\newfontfamily\cyrillicfont{Liberation Serif}
\newfontfamily\cyrillicfont{CMU Serif}	% Unicode Computer modern font
%Other fonts are too "fat"/wide

%% Liberation Serif for VeA title font
\definecolor{SeaBlue}{cmyk}{1,0,0.37,0.52}
%% Use proper Aurora Bold Condensed BT
\newfontfamily\TitleFont[%
	Path = /home/johnlm/Development/resources/tex_graphics_lib/fonts/ ,%
	FakeStretch = 1.5 ,
	Color = SeaBlue
	]{aurora-BT-condensed-bold}

% External common image library
\graphicspath{{/home/johnlm/Development/resources/tex_graphics_lib/}{img/}}

\hyphenpenalty=7500
\clubpenalty=9000
\widowpenalty=9000

% Set rubberband values
\setlength{\parskip}{1ex plus 3ex minus 1ex}	% Space between paragraphs


%\renewcommand{\thefootnote}{\alph{footnote}} % Piezīmes ar burtiem
\newcommand{\citeet}[1]{\textsuperscript{\cite{#1}}}

%% Numbering format fix for sections and subsections (trailing full-stop)
%\def\thesection{\arabic{section}.}
%\def\thesubsection{\thesection\arabic{subsection}.}
% Add deeper levels if required
%% FIXED: fixlatvian.sty provides these (REMOVE THIS)

% Reformatting captions
\captionsetup{format=hang,font={small}}	% GLOBAL
%\captionsetup[table]{position=top,

% Footer adjustements (by redefining plain pagestyle)
\setlength{\headheight}{15.2pt}
\fancypagestyle{plain}{%
	\fancyhf{}	% Clear current (default) header/footer
	\rfoot{\thepage}	% Page number on outside
	\renewcommand{\headrulewidth}{0pt}
	\renewcommand{\headrulewidth}{0pt}}


% Other stuff
\newcommand{\todo}{\textcolor{red}{TODO}}

%% Inline Word inclusion formatting
\newcommand{\termEn}[1]{\textenglish{\itshape {#1}}} % inline English

% Superscript cite
\makeatletter
	\renewcommand\@citess[1]{\textsuperscript{[#1]}}
\makeatother




%% This file contains special, document-specific preamble settings
% Binary code formatting for instruction table
\newcommand{\instr}[7]{\texttt{%
	\textcolor{purple}{#1}%		Group OPCODE
	\textcolor{blue}{#2}%		Vector OPCODE (part 1)
	\textcolor{cyan}{#3}%		Vector OPCODE (part 2)
	\textcolor{lightgray}{#4}%		Middle Slack bits
	\textcolor{OliveGreen}{#5}%		Argument 1
	\textcolor{Green}{#6}%	Argument 2
	\textcolor{lightgray}{#7}}}%		Trailing Slack bits
\newcommand{\mnem}[1]{\texttt{\bfseries #1}}

%% listings formateeshanas uzstaadiijumi
%% Stili
\lstset{%
	numbers=left, numberstyle=\tiny, numbersep=5pt,
	tabsize=2,%
	frame=lines,framerule=0.5pt,%
	belowcaptionskip=2ex,%
	basicstyle=\footnotesize\ttfamily,%
	lineskip=-0.5ex,
	keywordstyle={\bfseries\color{blue}},%
	keywordstyle=[2]{\color{Brown}},%
	keywordstyle=[3]{\color{OliveGreen}},%
	keywordstyle=[4]{\color{TealBlue}},%
	keywordstyle=[5]{\color{Sepia}},%
	emphstyle={\color{Bittersweet}},%
	emphstyle={[2]\color{PineGreen}},%
	emph={[3]{cpu_lib,bus_t,t_alu,t_reg,t_shift}},% My global emphs
	emphstyle={[3]\color{Green}},%
	commentstyle=\slshape\color{gray},
	stringstyle=\color{orange}} % "listings" koda stils
%% listings definiicijas
\lstdefinelanguage[qucs]{VHDL}[]{VHDL}%
{%morekeywords=[1]{},% MOAR core keywords
morekeywords=[3]{ieee,std_logic_1164,std_logic_arith,std_logic_unsigned,%
	work,std,textio,numeric_std,% usual libraries
	std_logic,std_logic_vector,unsigned,signed,%logic data types
	time,integer,real,line,string,positive,natural,% data types
	ns},%real type units
morekeywords=[4]{reverse_range,length,event,last_value,high,low},% attributes
morekeywords=[5]{write,writeline,conv_integer,to_unsigned,rising_edge,falling_edge}%,%funkcijas
%morestring=[b]',
}

\lstdefinelanguage{JSI}%
 {morekeywords={LD,LDI,MOV,ST,AR,ADD,SUB,INC,DEC,AND,OR,XOR,NOT,CLR,%
	MOV,LSL,LSR,ROR,ROL,BREQ,BRNQ,BRGT,BRGE,BRLE,BRLT,HLT,JMP},%
morekeywords=[2]{DEVREG_ADDR,SPI_RING_REGISTER,%
	BOOT_ROM_START,BOOT_ROM_END,SPECIAL_REG_START,SPECIAL_REG_END,%
	SRAM_START,SRAM_END},% Ports and internal registers
morekeywords=[3]{def,include,equ,dw,org},% directives
   sensitive=true,%
   %morecomment=[l]*,%
   morestring=[b]",
   morecomment=[l];%
   }[keywords,comments,strings]


%% Number within section
%\renewcommand{\thelstlisting}{\thesection.\arabic{lstlisting}}
\lstMakeShortInline[language={[qucs]VHDL},basicstyle=\normalsize\ttfamily]|

%% English determined number suffixes
%\renewcommand{\th}{\textsuperscript{th} }
\newcommand{\st}{\textsuperscript{st} }
\newcommand{\nd}{\textsuperscript{nd} }
\newcommand{\rd}{\textsuperscript{rd} }
\newcommand{\nth}{\textsuperscript{th} }

%% Abstract title style
\newcommand{\abstitlestyle}[1]{%
	\noindent \begin{center}
		\textbf{\Large #1}
	\end{center}}


% Hyperref package for links in PDF (should be last package)
\usepackage[linkbordercolor={blue},hyperindex]{hyperref}
\hypersetup{pdftitle={Iekļautās sistēmas mikrokontroliera kodola izstrāde}}
\hypersetup{pdfauthor={Jānis Šmēdiņš}}


% \onehalfspacing

% Chardump: „” em— en– fig‒ “”
\begin{document}
	% Titullapa
	\pagestyle{empty}
	% LM izcilā VeA titullapa bakalauram!
\begin{titlepage}
	% VeA Logo
	% "Ventspils Austskola" virsraksts (\TitleFont jābūt definētam!!!)
	\newsavebox{\veatext}
	\savebox{\veatext}{
		\TitleFont\Huge\MakeUppercase{Ventspils Augstskola}}
	% Teksta platuma noteikšana (lai pielīdzinātu bildi tā platumam}
	\newlength{\veatextwidth}
	\settowidth{\veatextwidth}{\usebox{\veatext}}
	% Šeit drukāts pats logo un nosaukums
	\centering\includegraphics[width=0.98\veatextwidth]{VeA_logo.pdf}\\[4pt]
	\usebox{\veatext}\\[6pt]
	\large Informācijas tehnoloģiju fakultāte\\[2cm]
	
	%\textbf{Maģistra darbs}\\[1.5cm]
	\textbf{Maģistra darba 1.~atskaite}\\[1.5cm]
	
	\textsc{\Large Attēlu raksturpunktu atrašanas un salāgošanas algoritmu veiktspējas izpēte uz dažādām aparatūras platformām}
	\vfill % LIELĀ atstarpe
	
	%\raggedleft
	\normalsize
	\begin{minipage}[t]{0.3\textwidth}
		\begin{flushleft}
			Autors
		\end{flushleft}
	\end{minipage}
	\begin{minipage}[t]{0.65\textwidth}
		\begin{flushleft}
			Ventspils Augstskolas \\
			Informācijas tehnoloģiju fakultātes \\
			profesionālās maģistra studiju programmas „Elektronika”\\
			2.~kursa students \\
			\textbf{Jānis Šmēdiņš}\\
			Matr.~nr.~\texttt{12140012}\\
			\rule[-1em]{10em}{1pt}\\
			\makebox[10em][c]{\tiny (paraksts)}\\[1cm]
		\end{flushleft}
	\end{minipage}\\[2em]
	\begin{minipage}[t]{0.3\textwidth}
		\begin{flushleft}
			Fakultātes dekāns
		\end{flushleft}
	\end{minipage}
	\begin{minipage}[t]{0.65\textwidth}
		\begin{flushleft}
			asoc.prof.,~Dr.~math.~Gaļina Hiļķeviča\\[1ex]
			\rule[-1em]{10em}{1pt}\\
			\makebox[10em][c]{\tiny (paraksts)}\\[1cm]
		\end{flushleft}
	\end{minipage}\\[2em]
	\begin{minipage}[t]{0.3\textwidth}
		\begin{flushleft}
			Zinātniskais vadītājs
		\end{flushleft}
	\end{minipage}
	\begin{minipage}[t]{0.65\textwidth}
		\begin{flushleft}
			pētnieks, Dr.~sc.~comp.~Kaspars Sudars\\[1ex]
			\rule[-1em]{10em}{1pt}\\
			\makebox[10em][c]{\tiny (paraksts)}\\[1cm]
		\end{flushleft}
	\end{minipage}\\[2em]
	\begin{minipage}[t]{0.3\textwidth}
		\begin{flushleft}
			Recenzents
		\end{flushleft}
	\end{minipage}
	\begin{minipage}[t]{0.65\textwidth}
		\begin{flushleft}
			\rule[-1em]{20em}{1pt}\\
			\makebox[20em][c]{\tiny
				(ieņemamais amats,
				zinātniskais nosaukums,
				vārds, uzvārds)}\\[1ex]
			\rule[-1em]{10em}{1pt}\\
			\makebox[10em][c]{\tiny (paraksts)}\\[1cm]
		\end{flushleft}
	\end{minipage}\\[1cm]
	
	\centering
	Ventspils\\
	\the\year
\end{titlepage}
 % Ielikt titullapu
	\stepcounter{page} %Palielināt lapu skaitītāju (jeb ieskaitīt titullapu)
	\lstset{%
			numbers=left, numberstyle=\tiny, numbersep=5pt,
			tabsize=4,%
			basicstyle=\footnotesize\ttfamily,%
			keywordstyle={\bfseries\color{blue}},%
			keywordstyle=[2]{\color{Brown}},%
			keywordstyle=[3]{\color{OliveGreen}},%
			keywordstyle=[4]{\color{Bittersweet}},%
			commentstyle=\slshape\color{gray},%
			emphstyle={\color{PineGreen}},%
			stringstyle=\color{orange}} % "listings" koda stils
	
	%\section*{Saturs}
	\renewcommand{\contentsname}%
		{\vspace*{-12mm}\section*{Saturs}\vspace*{-6mm}}
	\tableofcontents
	%\addcontentsline{toc}{section}{Saturs}	% Saturs ir saturaa?
	
	\clearpage
	%\sloppy %Don't care too much about rigid formatting
	\pagestyle{plain}
	\onehalfspacing % 1.5 spacing
	
	% Termini
		% PCI, PCI Express
		% Seriālā, paralēlā datu pāraide
	
	%\section*{Anotācija}
	%\renewcommand{\abstractname}%
	%	{\section*{Anotācija}}
	%\begin{abstract}
	\addcontentsline{toc}{section}{Anotācija}
	\todo
\end{abstract}

\clearpage
\begin{english}
	\begin{abstract}
		\addcontentsline{toc}{section}{Abstract}
		\todo
	\end{abstract}
\end{english}

\clearpage
\begin{russian}
	\begin{abstract}
		\addcontentsline{toc}{section}{Аннотация}
		\todo \\
		Чо такоё?
	\end{abstract}
\end{russian}

	
	\clearpage
	%\section*{Ievads} \addcontentsline{toc}{section}{Ievads}
% FIXME: Should I do the "screaming" first sentence?
Zinātne nestāv uz vietas --- nepārtraukti tiek uzsākti jauni zinātniskie projekti,
un arvien biežāk tajos nepieciešami specializēti elektroniskie risinājumi.
Šie risinājumi ir projekta specifiski, bet tajos bieži ir nepieciešamas
kontroles iekārtas, datu formēšanas un pārraides iekārtas, un
šim pielietojumam var izmantot specializētus mikrokontrolierus.
Šādu mikrokontrolieru izstrāde tiek krietni vienkāršota, ja pieejams
modulārs mikrokontroliera kodols ar vienkāršu un adaptējamu saskarni.

Sevišķas prasības specializētu mikrokontrolieru izstrādei ir kosmosa
tehnoloģijās, kur jaunas ierīces un komponentes pieprasa īpaši rūpīgu
pārbaudi, uzticama mikrokontroliera kodola universalitāte ir sevišķi
vērtīga.

Šī bakalaura darba mērķis ir, pirmkārt, izstrādāt modulāru mikrokontroliera kodolu,
kurš, bez vai ar minimālām modifikācijām, būtu izmantojams,
vispārēja un specializēta pielietojuma, mikrokontrolieru implemetācijās,
sintēzei FPGA (\termEn{{field-programmable} gate array}), un otrkārt,
izstrādāt vienkāršu parauga mikrokontrolieri, kas šo kodolu izmanto,
demonstrējot izstrādātā kodola izmantošanu šim nolūkam.
Mērķa sasniegšanai ir izvirzīti vairāki uzdevumi.
\begin{enumerate}
	\item Kodola un perifērijas saskarnes definēšana,
		kurai jābūt pietiekami vienkāršai, lai atbalstītu plašu spektru
		komponenšu.
	\item Kodola izmantojamās instrukciju kopas definēšana. Tai jābūt
		vienkāršai un jānodrošina pamata funkcionalitāti.
	\item Kodola arhitektūras definēšana un izstrāde. Tas galvenokārt nozīmē
		visu kodola apakškomponenšu izstrādi un komplektēšanu.
	\item Parauga mikrokontroliera uzbūves definēšana un komponenšu izstrāde.
\end{enumerate}

Darba izstrādes galvenais instruments ir VHDL, kas ir
aparatūras apraksta valoda (HDL) ar kuru var aprakstīt shēmas darbību un uzbūvi.
HDL nozīmīgākā īpašība ir shēmas aprakstu sintezējamība fiziskā realizācijā,
galvenokārt FPGA mikroshēmā, kas ir darba mērķa platforma.
FPGA ir integrētā shēma, kuru vienkāršoti var 
uzskatīt par ,,tukšu'', ar HDL palīdzību, programmējamu mikroshēmu.

Darbā apskatīta eksistējoša, HDL aprakstīta procesora prototipa arhitektūra,
veikta tās analīze, ieviešot arhitektūras optimizācijas un risinājumus tās
trūkumiem. Šī prototips, ar izvirzītajiem arhitektūras modifikācijām un 
uzlabojumiem, izmantots par bāzi izstrādājamā kodola izstrādē.

%Šī darba izklāsta sākumā tiek apskatītas HDL, to pielietojums un divu
%populārāko valodu --- VHDL un Verilog --- salīdzinājums.
%Darba turpinājumā tiek apskatīta HDL sintēzes mērķa platformu ---
%FPGA un ASIC --- uzbūve un tās iespaidu uz izstrādi.

Šis darbs ir sadalāms divās daļās, no kurām pirmā sniedz ieskatu
izmantotajās tehnoloģijās --- HDL, to pielietojumu un uzbūvi, kā arī HDL
sintēzes mērķa platformu (t.sk.,~FPGA) uzbūvi. Otrā daļā tiek veikta bāzes
prototipa analīze un izvirzīti uzlabojumi, pēc kuras izklāstīta izstrādātā
kodola un parauga mikrokontroliera implementācija.

 % TODO: Rewrite this bugger
	
	\clearpage
	\section{Sistēmas perifērija}
	Atsevišķs procesors ir nefunkcionāls, ja tam nevar pievadīt izpildāmās
	instrukcijas, atgūt datus vai vadīt kādu citu perifēro ierīci.
	
	Praktiski vienmēr procesoram nepieciešama atmiņas iekārta (RAM), kur
	glabāt izpildāmo programmu, kā arī pirmsapstrādes, pēcapstrādes un 
	apstrādes laika datus.
	
	Šajā nodaļā apskatītas divas sistēmas shēmas, no kurām viena veiksmīgi
	nosimulēta, bet nav sitezējama (rev.~02) un krietni sarežģītākā
	sintezējamā shēma, kura satur papildus perifērās ierīces.
	
	\subsection{Simulētā shēma (rev.~02)}
Sākotnējā darba augšējā līmeņa sistēmas gala shēma paredzēta, kā
paškomplektējoša, tikai ar procesoru un atmiņu, 
kur no ārpuses tiek tikai dota takts (\texttt{CLOCK})
un atiestatīšanas signāls ($\overline{\texttt{RESET}}$). Darbības pārbaudei
tiek izmantotas izstrādes platformas piedāvātās diodes.

\begin{figure}[bhp]
	\centering
	\def\svgwidth{\textwidth}
	{\ttfamily\tiny\input{img/top-rev2.pdf_tex}}
	\caption{Augšējā līmeņa shēma (rev.~02).}
	\label{fig:top-rev2}
\end{figure}

Šī shēma tika veiksmīgi simulēta pirms sintēzes
(rezultātus sk.~\ref{appx:simulation}.~pielikumā),
bet tā nav korekti sintezējama, jo paredz RAM sākotnējos datus,
kurus sintēzes rīks ignorē.
Realitātē RAM dati tiek pazaudēti tikko tiek noņemts barošanas spriegums un
tātad pēc šādas shēmas nav iespējams saglabāt izpildāmo programmu.

\pagebreak[3]
%\subsubsection{Operatīvā atmiņa}
	Operatīvā atmiņa šeit realizēta ar divām vienvirziena datu apmaiņas
	šinām. Kontroles signāli izmantoti līdzīgi klasiskai trīs-stāvokļu
	divvirzienu datu šinas atmiņai, pielāgojoties procesora atmiņas
	saskarnei.
	
	\begin{singlespace}
		\lstinputlisting[language={[qucs]VHDL},%float=p,
		                caption={RAM VHDL entītija.},%
		                label=kb:ram-entity,%
		                linerange={7-13},firstnumber=7,
		                breaklines,breakatwhitespace]
			{code/mem.twoport.vhd}
	
		\lstinputlisting[language={[qucs]VHDL},%float=p,
		                caption={RAM VHDL arhitektūras apraksts (izgriezums).},%
		                label=kb:ram-trimmed,%
		                linerange={85-99},firstnumber=85,%
		                breaklines,breakatwhitespace]
			{code/mem.twoport.vhd}
	\end{singlespace}
	
	\noindent Pilno kodu skatīt \ref{appx:ram-code}.~pielikumā.
 \clearpage %\pagebreak[3]
	\subsection{Papildinātā shēma (rev.~03)}
Šī apakšnodaļa apraksta sistēmas papildinājumus funkcionālas, sintezējamas
sistēmas realizācijai. Šīs revīzijas implementācija \textbf{nav pabeigta},
tādēļ šī apakšnodaļa ir uzskatāma par \textbf{turpmākā darba dokumentāciju},
nevis realizētu implementāciju.

Trešā sistēmas revīzija pievieno patstāvīgās atmiņas elementus sākotnējās
programmas un datu uzglabāšanai. Šīs atmiņas ierīces piekārtotas kopējai
adresācijas telpai, tādējādi rodas nepieciešamība pēc adreses dekodēšanas
loģikas, kuru nodrošina MMU (\termEn{Memory map unit}).

\begin{figure}[bhp]
	\centering
	\def\svgscale{1.25}
	{\input{img/top-rev3.pdf_tex}}
	\caption{Augšējā līmeņa blokshēma (rev.~03).}
	\label{fig:top-rev3}
\end{figure}

\subsubsection{Boot ROM}
	\termEn{Boot ROM} ir tikai nolasāmā atmiņa, kura paredzēta
	\termEn{bootstrap} procesa programmas uzglabāšanai.
	Šīs programmas uzdevums ir ielādēt reālo izpildes programmu no
	patstāvīgās atmiņas — šajā implementācijā konkrēti 
	no SPI \termEn{Flash} atmiņas.
	
	Par \termEn{Boot ROM} kalpos \termEn{Actel Fusion} piedāvātais
	\termEn{FlashROM} makross.\citeet{FlashROM}
	Maksimāli pieejamais ROM garums uz izmantojamās platformas ir
	1 kilobits jeb 64 vārdi.\citeet{FusionGuide}

\subsubsection{Adrešu telpai piekārtotā ievade/izvade}
	Papildus adresācijas telpā tiks piekārtoti speciālie reģistri dažādu
	perifēro ierīču datu apmaiņai un konfigurācijai.
	
	Uz doto brīdi (un kā redzams blokshēmā — \ref{fig:top-rev3}.~att.)
	paredzēta tikai SPI saskarne. Tam paredzēts viens pārraides un viens
	konfigurācijas reģistrs. Tam paredzēta iespēja nokonfigurēt datu
	apmaiņu pa baitam (8 biti) vai pa vārdam (16 biti). Datu pārraide tiek
	uzsākta uzreiz pēc jauno datu ierakstes reģistrā (ja SPI iespējots) un
	tiek pārraidīti bez procesora līdzdalības.
	
	\begin{figure}[thp]
		\centering
		%\def\svgwidth{7cm}
		\def\svgscale{1.25}
		{\ttfamily\scriptsize\input{img/sub-spi.pdf_tex}}
		\caption{SPI saskarnes ierīce.}
		\label{fig:spi}
	\end{figure}
	
	SPI saskarne nepieciešama komunikācijai ar SPI \termEn{Flash} atmiņu,
	kura pieejama uz izstrādes platformas,\cite[43.~lpp.]{FusionGuide}
	kur tiks glabāta galvenā izpildāmā programma, kura \termEn{boot} procesa
	laikā tiks pārvietota uz RAM.

\subsubsection{MMU — adrešu dekoderis}
	MMU jeb adrešu dekoderis ir ierīce, kas atbild par adrešu telpas
	pārdalīšanu dažādām ierīcēm, kā arī nodrošināt korektu komunikāciju
	ar šīm ierīcēm.
	
	\begin{figure}[thp]
		\centering
		\def\svgwidth{0.75\textwidth}
		{\ttfamily\small\input{img/remap.pdf_tex}}
		\caption{Adrešu telpas sadalījums.}
		\label{fig:memory-map}
	\end{figure}
	
	MMU tādējādi ar procesoru veido tādu pašu saskarni kā RAM un aparatūras
	līmenī ir pilnībā „caursīdīga”, savukārt piekārtojamo ierīču saskarne
	ir implementācijas definēta.
	
	\termEn{Boot ROM} nepieciešams piekārtot ar sākumu adresē \texttt{0x0000},
	lai pēc atiestatīšanas procesors sāktu izpildīt \termEn{boot} programmu.
 \clearpage %\pagebreak[3]
	\subsection{Revīziju vēsture} \label{sec:sys-revs}
%\todo

\begin{description}
	\item[Rev.~01] \hfill \\
		Kā RAM izmantota atmiņas realizācija ar 
		trīs-stāvokļu, divvirzienu datu šinu.
	\item[Rev.~02] \hfill \\
		Līdzīgi revīzijas izmaiņām procesora uzbūvē, arī atmiņa pārveidota,
		likvidējot trīs-stāvokļu šinu un tā vietā divas vienvirziena šinas
		datu ieejai un izejai. Kontroles signāli atstāti bez izmaiņām,
		tādējādi atmiņa vēljoprojām „emulē” trīs-stāvokļu datu apmaiņu,
		saglabājot procesora iespēju strādāt ar divvirzienu šinas atmiņu
		(papildus izmantojot virziena maiņas buferus).
	\item[Rev.~03] \hfill \\
		Šī revīzija veic fundamentālas izmaiņas sistēmas perifērijā,
		pievienojot patstāvīgo atmiņu un realizējot adrešu telpas
		piekārtošanu datu ievades/izvades ierīcēm ar MMU palīdzību.
		Šai arī paredz sistēmas \termEn{boot} procesu, kur izpildes
		programma tiek pārvietota izpildei uz operatīvo atmiņu no
		patstāvīgās atmiņas.
\end{description}

	
	% Literatūras saraksts
	\bibliographystyle{ieeetr}
	%\renewcommand{\refname}%
	%	{\vspace*{-12mm}\section*{Izmantotā literatūra}\vspace*{-6mm}}
	\clearpage
	%\pagestyle{empty}
	\raggedright
	\begin{thebibliography}{9}
		\addcontentsline{toc}{section}{\refname}
		\bibitem{Perry-VHDL}
			Douglas L.~Perry,
			\textit{VHDL: Programming by Example}, 4\nth Edition. \linebreak[2]
			New York: McGraw-Hill, 2002. ISBN~0-07-140070-2
		
		\bibitem{FusionGuide}
			Actel corp.,
			\textit{Fusion Embedded Development Kit User's Guide}. %\linebreak[2]
			%USA: Actel,
			2009.
		
		\bibitem{FusionFAQ}
			Actel corp.,
			\textit{Synplify — Synthesis Frequently Asked Questions}. %\linebreak[2]
			%USA: Actel,
			2009.
		
		\bibitem{FlashROM}
			Actel corp.,
			\textit{Fusion FlashROM}, Application Note AC236. %\linebreak[2]
			%USA: Actel,
			2005.
		
		%\bibitem{LDD3}
		%	Corbet~J., Rubini~A., Kroah-Hartman~G.,
		%	\textit{Linux Device Drivers, Third Edition}.
		%	Sebastopol: O'Reilly, 2005. ISBN~0-596-00590-3
		
		%\bibitem{PCIeArch}
		%	Burduk~R., Anderson~D., Shanley~T.,
		%	\textit{PCI Express System Architecture}.
		%	\linebreak[3]
		%	Boston: Addison-Wesley, 2009. ISBN~0-321-15630-7
	\end{thebibliography}
	
\end{document}
